\ifdefined\gkver\else\def\gkver{student}\fi
\documentclass[\gkver]{gaokaozhenti}
\begin{document}

\chapter{1993年全国}

\begin{enumerate}
\item
%% number: 1
%% paperName: 1993年全国高考第 1 题 2 分
%% typeId: 1
%% type: 单选题
%% chapter: 磁场
%% point: 带电粒子在磁场中的运动
%% method: 无
%% score: 2
%% degree: 650
%% duplicateId: 0
%% body:
两个电子以大小不同的初速度沿垂直于磁场的方向射入同一匀强磁场中。设 $r_{1}$、$r_{2}$ 为这两个电子的运动轨道半径，$T_{1}$、$T_{2}$ 是它们的运动周期，则\xzanswer{D}
\fourchoices[answer=D]
{$r_{1}=r_{2}$，$T_{1}\neq T_{2}$}
{$r_{1}\neq r_{2}$，$T_{1}\neq T_{2}$}
{$r_{1}=r_{2}$，$T_{1}=T_{2}$}
{$r_{1}\neq r_{2}$，$T_{1}=T_{2}$}

%% answer: D
%% memo:
\memoanswer{
由 $qvB=m\frac{v^{2}}{r}$ 得 $r=\frac{mv}{qB}$，因 $v_{1}\neq v_{2}$，得 $r_{1}\neq r_{2}$；又 $T=\frac{2\pi m}{qB}$，得 $T_{1}=T_{2}$，故 D 正确。
}

\item
%% number: 2
%% paperName: 1993年全国高考第 2 题 2 分
%% typeId: 1
%% type: 单选题
%% chapter: 曲线运动
%% point: 人造地球卫星
%% method: 无
%% score: 2
%% degree: 650
%% duplicateId: 0
%% body:
同步卫星是指相对于地面不动的人造地球卫星\xzanswer{D}
\fourchoices[answer=D]
{它可以在地面上任一点的正上方，且离地心的距离可按需要选择不同值}
{它可以在地面上任一点的正上方，但离地心的距离是一定的}
{它只能在赤道的正上方，但离地心的距离可按需要选择不同值}
{它只能在赤道的正上方，且离地心的距离是一定的}

%% answer: D
%% memo:
\memoanswer{
由于同步卫星所受地球的万有引力作为向心力，故卫星必须在赤道所在的圆平面内，因此在赤道正上方，且 $G\frac{Mm}{r^{2}}=m\frac{4\pi^{2}}{T^{2}}r$，所以 $r=\sqrt[3]{\frac{GMT^{2}}{4\pi^{2}}}$，由于 $T=24$ 小时为定值，故 $r$ 为定值，计算得 $r=4.2\times10^{7}\Um$，故 D 正确。
}

\item
%% number: 3
%% paperName: 1993年全国高考第 3 题 2 分
%% typeId: 1
%% type: 单选题
%% chapter: 交流电
%% point: LC振荡回路
%% method: 无
%% score: 2
%% degree: 650
%% duplicateId: 0
%% body:
由自感系数为 $L$ 的线圈和可变电容器 $C$ 构成收音机的调谐电路。为使收音机能接收到 $f_{1}=550\UkHz$ 至 $f_{2}=1650\UkHz$ 范围内的所有电台的播音，则可变电容器与 $f_{1}$ 对应的电容 $C_{1}$ 和与 $f_{2}$ 对应的电容 $C_{2}$ 之比为\xzanswer{D}
\fourchoices[answer=D]
{$1:\sqrt{3}$}
{$\sqrt{3}:1$}
{$1:9$}
{$9:1$}

%% answer: D
%% memo:
\memoanswer{
LC 振荡电路频率为 $f=\frac{1}{2\pi\sqrt{LC}}$，得 $C=\frac{1}{4\pi^{2}f^{2}L}$，所以 $\frac{C_{1}}{C_{2}}=\left(\frac{f_{2}}{f_{1}}\right)^{2}=\left(\frac{1650}{550}\right)^{2}=\frac{9}{1}$，故 D 正确。
}

\item
%% number: 4
%% paperName: 1993年全国高考第 4 题 2 分
%% typeId: 1
%% type: 单选题
%% chapter: 原子和原子核
%% point: 半衰期
%% method: 无
%% score: 2
%% degree: 650
%% duplicateId: 0
%% body:
若元素 A 的半衰期为 4 天，元素 B 的半衰期为 5 天，则相同质量的 A 和 B，经过 20 天后，剩下的质量之比 $m_{A}:m_{B}$ 为\xzanswer{C}
\fourchoices[answer=C]
{$30:31$}
{$31:30$}
{$1:2$}
{$2:1$}

%% answer: C
%% memo:
\memoanswer{
由原子核的衰变公式 $m=m_{0}\left(\frac{1}{2}\right)^{\frac{t}{T}}$，所以 $\frac{m_{A}}{m_{B}}=\frac{\left(\frac{1}{2}\right)^{\frac{20}{4}}}{\left(\frac{1}{2}\right)^{\frac{20}{5}}}=\frac{1}{2}$，故 C 正确。
}

\item
%% number: 5
%% paperName: 1993年全国高考第 5 题 2 分
%% typeId: 1
%% type: 单选题
%% chapter: 光学
%% point: 光的干涉
%% method: 无
%% score: 2
%% degree: 750
%% duplicateId: 0
%% body:
图中所示是用干涉法检查某块厚玻璃板的上表面是否平的装置。所用单色光是用普通光源加滤光片产生的。检查中所观察到的干涉条纹是由下列哪两个表面反射的光线叠加而成的？\xzanswer{C}
\onepicture[w=5.5cm]{\includesvg{figs/05}}
\fourchoices[answer=C]
{a 的上表面和 b 的下表面}
{a 的上表面和 b 的上表面}
{a 的下表面和 b 的上表面}
{a 的下表面和 b 的下表面}

%% answer: C
%% memo:
\memoanswer{
由薄膜干涉原理知 C 正确，因为这两个表面形成厚度均匀变化的空气层。
}

\item
%% number: 6
%% paperName: 1993年全国高考第 6 题 2 分
%% typeId: 1
%% type: 单选题
%% chapter: 光学
%% point: 透镜
%% method: 无
%% score: 2
%% degree: 750
%% duplicateId: 0
%% body:
一物体经凸透镜在屏上成一放大的实像。凸透镜主轴沿水平方向。今将凸透镜向上移动少许，则\xzanswer{A}
\fourchoices[answer=A]
{屏上像的位置向上移动}
{屏上像的位置向下移动}
{屏上像的位置保持不动，但像变大}
{屏上像的位置保持不动，但像变小}

%% answer: A
%% memo:
\memoanswer{
可利用作图法求出屏上像的位置向上移动，故 A 正确。
}

\item
%% number: 7
%% paperName: 1993年全国高考第 7 题 2 分
%% typeId: 1
%% type: 单选题
%% chapter: 磁场
%% point: 安培力
%% method: 无
%% score: 2
%% degree: 650
%% duplicateId: 0
%% body:
如图所示的电流天平可用来测定磁感应强度 $B$。天平的右臂下面挂有一个矩形线圈，宽 ab 为 $l$，共 $N$ 匝，线圈的下部悬在待测匀强磁场中，磁场方向垂直纸面。当线圈中通有图示方向电流 $I$ 时，在天平左、右两边加上质量各为 $m_{1}$、$m_{2}$ 的砝码，天平平衡。当电流反向（大小不变）时，右边再加上质量为 $m$ 的砝码后，天平重新平衡。由此可知\xzanswer{B}
\onepicture[w=5.5cm]{figs/07.png}
\fourchoices[answer=B]
{磁感应强度的方向垂直纸面向里，大小为 $\frac{(m_{1}-m_{2})g}{NIl}$}
{磁感应强度的方向垂直纸面向里，大小为 $\frac{mg}{2NIl}$}
{磁感应强度的方向垂直纸面向外，大小为 $\frac{(m_{1}-m_{2})g}{NIl}$}
{磁感应强度的方向垂直纸面向外，大小为 $\frac{mg}{2NIl}$}

%% answer: B
%% memo:
\memoanswer{
由题意知线圈下边受安培力方向在电流方向改变后向上。用左手定则判定磁感应强度垂直纸面向里，且 $mg=2NBIl$，解得 $B=\frac{mg}{2NIl}$，故 B 正确。
}

\item
%% number: 8
%% paperName: 1993年全国高考第 8 题 2 分
%% typeId: 1
%% type: 单选题
%% chapter: 机械波
%% point: 波的图象
%% method: 无
%% score: 2
%% degree: 750
%% duplicateId: 0
%% body:
一列沿 $x$ 方向传播的横波，其振幅为 $A$，波长为 $\lambda$，某一时刻波的图象如图所示。在该时刻，某一质点的坐标为（$\lambda$，0），经过四分之一周期后，该质点的坐标为\xzanswer{B}
\onepicture[w=6cm]{\includesvg{figs/08}}
\fourchoices[answer=B]
{$\frac{5}{4}\lambda$，0}
{$\lambda$，$-A$}
{$\lambda$，$A$}
{$\frac{5}{4}\lambda$，$A$}

%% answer: B
%% memo:
\memoanswer{
方法一：研究波形——经四分之一周期后，波向右移动了四分之一波长，如图所示，A 点的坐标为（$\lambda$，$-A$）。

\onepicture[w=6cm]{\includesvg{figs/08_an}}

方法二：研究质点——由微平移法可得，此刻 A 点向下振动，经过四分之一周期后，达到波谷，横坐标不变，故 A 点的坐标为（$\lambda$，$-A$）。
}

\item
%% number: 9
%% paperName: 1993年全国高考第 9 题 2 分
%% typeId: 1
%% type: 单选题
%% chapter: 电路
%% point: 多用表的使用
%% method: 无
%% score: 2
%% degree: 550
%% duplicateId: 0
%% body:
下图为万用表欧姆挡的原理示意图，其中电流表的满偏电流为 $300\UmA$，内阻 $r_{g}=100\UO$，调零电阻最大阻值 $R=50\UkO$，串联的固定电阻 $R_{0}=50\UO$，电池电动势 $\varepsilon=1.5\UV$。用它测量电阻 $R_{x}$，能准确测量的阻值范围是\xzanswer{B}
\onepicture[w=5cm]{figs/09.png}
\fourchoices[answer=B]
{$30\UkO\sim80\UkO$}
{$3\UkO\sim8\UkO$}
{$300\UO\sim800\UO$}
{$30\UO\sim80\UO$}

%% answer: B
%% memo:
\memoanswer{
该欧姆表的中值电阻即为欧姆表的总电阻，亦即为满偏时的总内阻，有 $R=\frac{E}{I_{g}}=\frac{1.5}{300\times10^{-6}}\UO=5000\UO$，故该表能准确测量的阻值范围包含 $5\UkO$，故 B 正确。
}

\item
%% number: 10
%% paperName: 1993年全国高考第 10 题 2 分
%% typeId: 1
%% type: 单选题
%% chapter: 力物体的平衡
%% point: 物体系平衡
%% method: 无
%% score: 2
%% degree: 550
%% duplicateId: 0
%% body:
A、B、C 三物块质量分别为 $M$、$m$ 和 $m_{0}$，作如图所示的联结。绳子不可伸长，且绳子和滑轮的质量、滑轮的摩擦均可不计。若 B 随 A 一起沿水平桌面作匀速运动，则可以断定\xzanswer{A}
\onepicture[w=5.5cm]{figs/10.png}
\fourchoices[answer=A]
{物块 A 与桌面之间有摩擦力，大小为 $m_{0}g$}
{物块 A 与 B 之间有摩擦力，大小为 $m_{0}g$}
{桌面对 A，B 对 A，都有摩擦力，两者方向相同，合力为 $m_{0}g$}
{桌面对 A，B 对 A，都有摩擦力，两者方向相反，合力为 $m_{0}g$}

%% answer: A
%% memo:
\memoanswer{
因 A 做匀速运动，受绳拉力为 $m_{0}g$，桌面对它的摩擦力大小应为 $m_{0}g$；B 做匀速运动，故不受 A 的摩擦力，靠惯性运动，否则没有其他力与此摩擦力平衡，故 A 正确。
}

\item
%% number: 11
%% paperName: 1993年全国高考第 11 题 2 分
%% typeId: 1
%% type: 单选题
%% chapter: 电场
%% point: 电场叠加
%% method: 无
%% score: 2
%% degree: 550
%% duplicateId: 0
%% body:
图中接地金属球 A 的半径为 $R$，球外点电荷的电量为 $Q$，到球心的距离为 $r$。该点电荷的电场在球心的场强等于\xzanswer{D}
\onepicture[w=5.5cm]{figs/11.png}
\fourchoices[answer=D]
{$\frac{kQ}{r^{2}}-\frac{kQ}{R^{2}}$}
{$\frac{kQ}{r^{2}}+\frac{kQ}{R^{2}}$}
{0}
{$\frac{kQ}{r^{2}}$}

%% answer: D
%% memo:
\memoanswer{
由静电平衡知，球心的合场强为 0，即金属球在 O 点产生的场强与点电荷 $Q$ 在 O 点产生的场强等大反向。由点电荷的场强公式知 $E=k\frac{Q}{r^{2}}$，故 D 正确。
}

\item
%% number: 12
%% paperName: 1993年全国高考第 12 题 2 分
%% typeId: 1
%% type: 单选题
%% chapter: 功和能
%% point: 功
%% method: 无
%% score: 2
%% degree: 650
%% duplicateId: 0
%% body:
小物块位于光滑的斜面上，斜面位于光滑的水平地面上。从地面上看，在小物块沿斜面下滑的过程中，斜面对小物块的作用力\xzanswer{B}
\onepicture[w=4.5cm]{\includesvg{figs/12}}
\fourchoices[answer=B]
{垂直于接触面，做功为零}
{垂直于接触面，做功不为零}
{不垂直于接触面，做功为零}
{不垂直于接触面，做功不为零}

%% answer: B
%% memo:
\memoanswer{
因斜面光滑，小物块不受摩擦力，故斜面对小物块的作用力即为支持力，应垂直于接触面。由于水平地面光滑，小物块下滑过程中，斜面向右运动（水平方向动量守恒），则作用力不垂直于木块的位移，做功不为零，故 B 正确。
}

\item
%% number: 13
%% paperName: 1993年全国高考第 13 题 2 分
%% typeId: 1
%% type: 单选题
%% chapter: 气体性质
%% point: 热力学第一定律
%% method: 无
%% score: 2
%% degree: 450
%% duplicateId: 0
%% body:
图中容器 A、B 各有一个可自由移动的轻活塞，活塞下面是水，上面是大气，大气压恒定。A、B 的底部由带有阀门 K 的管道相连。整个装置与外界绝热。原先，A 中水面比 B 中的高。打开阀门，使 A 中的水逐渐向 B 中流，最后达到平衡。在这个过程中，\xzanswer{D}
\onepicture[w=5.5cm]{figs/13.png}
\fourchoices[answer=D]
{大气压力对水做功，水的内能增加}
{水克服大气压力做功，水的内能减少}
{大气压力对水不做功，水的内能不变}
{大气压力对水不做功，水的内能增加}

%% answer: D
%% memo:
\memoanswer{
由题意知，大气对两活塞产生的压强为 $p_{A}=p_{B}=p$，产生的压力为 $F_{A}=pS_{A}$、$F_{B}=pS_{B}$，又因液体体积一定，故 $h_{A}=\frac{V}{S_{A}}$、$h_{B}=\frac{V}{S_{B}}$，所以 $W_{A}=F_{A}h_{A}$、$W_{B}=F_{B}h_{B}$，联立得 $W_{A}=W_{B}$，即大气压力对水做的正功、负功相等，总体来说不做功。由于达到平衡后水的质心的高度降低，故水的势能减少，由能量守恒知水的内能应增加，故 D 正确。
}

\item
%% number: 14
%% paperName: 1993年全国高考第 14 题 4 分
%% typeId: 1
%% type: 单选题
%% chapter: 光学
%% point: 光电效应
%% method: 无
%% score: 4
%% degree: 650
%% duplicateId: 0
%% body:
入射光照射到某金属表面上发生光电效应，若入射光的强度减弱，而频率保持不变，那么\xzanswer{C}
\fourchoices[answer=C]
{从光照至金属表面上到发射出光电子之间的时间间隔将明显增加}
{逸出的光电子的最大初动能将减小}
{单位时间内从金属表面逸出的光电子数目将减少}
{有可能不发生光电效应}

%% answer: C
%% memo:
\memoanswer{
入射光的强度减弱，则单位时间内入射的光子数减少，由光电效应的规律知，单位时间内从金属表面逸出的光电子数目将减少，故 C 正确、D 错误；光电效应是瞬时发生的，故 A 错误；光电子的最大初动能与入射光的频率有关，与强度无关，故 B 错误。
}

\item
%% number: 15
%% paperName: 1993年全国高考第 15 题 4 分
%% typeId: 1
%% type: 单选题
%% chapter: 气体性质
%% point: 分子力
%% method: 无
%% score: 4
%% degree: 650
%% duplicateId: 0
%% body:
分子间的相互作用力由引力 $f_{\text{引}}$ 和斥力 $f_{\text{斥}}$ 两部分组成，则\xzanswer{A}
\fourchoices[answer=A]
{$f_{\text{斥}}$ 和 $f_{\text{引}}$ 是同时存在的}
{$f_{\text{引}}$ 总是大于 $f_{\text{斥}}$，其合力总表现为引力}
{分子之间的距离越小，$f_{\text{引}}$ 越小，$f_{\text{斥}}$ 越大}
{分子之间的距离越小，$f_{\text{引}}$ 越大，$f_{\text{斥}}$ 越小}

%% answer: A
%% memo:
\memoanswer{
分子间斥力 $F_{\text{斥}}$ 与引力 $F_{\text{引}}$ 同时存在，存在一个距离 $r_{0}$，当 $r=r_{0}$ 时，$F_{\text{斥}}$ 与 $F_{\text{引}}$ 相等，分子力表现为零；当 $r>r_{0}$ 时，分子间距离越大，$F_{\text{斥}}$ 与 $F_{\text{引}}$ 都减小，但 $F_{\text{斥}}$ 减小得快，故表现为引力；当 $r<r_{0}$ 时，分子间距离越小，$F_{\text{斥}}$ 与 $F_{\text{引}}$ 都增大，但 $F_{\text{斥}}$ 增大得快，表现为斥力，故 A 正确。
}

\item
%% number: 16
%% paperName: 1993年全国高考第 16 题 4 分
%% typeId: 2
%% type: 多选题
%% chapter: 交流电
%% point: 变压器
%% method: 无
%% score: 4
%% degree: 550
%% duplicateId: 0
%% body:
如图所示，一理想变压器的原、副线圈分别由双线圈 ab 和 cd（匝数都为 $n_{1}$）、ef 和 gh（匝数都为 $n_{2}$）组成。用 $I_{1}$ 和 $U_{1}$ 表示输入电流和电压，$I_{2}$ 和 $U_{2}$ 表示输出电流和电压。在下列四种连接法中，符合关系 $\frac{U_{1}}{U_{2}}=\frac{n_{1}}{n_{2}}$、$\frac{I_{1}}{I_{2}}=\frac{n_{2}}{n_{1}}$ 的有\xzanswer{AD}
\onepicture[w=5cm]{figs/16.png}
\fourchoices[answer=AD]
{b 与 c 相连，以 a、d 为输入端；f 与 g 相连，以 e、h 为输出端}
{b 与 c 相连，以 a、d 为输入端；e 与 g 相连、f 与 h 相连作为输出端}
{a 与 c 相连、b 与 d 相连作为输入端；f 与 g 相连，以 e、h 为输出端}
{a 与 c 相连、b 与 d 相连作为输入端；e 与 g 相连、f 与 h 相连作为输出端}

%% answer: AD
%% memo:
\memoanswer{
用逆向思维解题，A 项：初级匝数为 $2n_{1}$，次级为 $2n_{2}$；B 项：初级为 $2n_{1}$，次级为 $n_{2}$；C 项：初级为 $n_{1}$，次级为 $2n_{2}$；D 项：初级为 $n_{1}$，次级为 $n_{2}$，故 AD 正确。
}

\item
%% number: 17
%% paperName: 1993年全国高考第 17 题 4 分
%% typeId: 2
%% type: 多选题
%% chapter: 电路
%% point: 电路中的图象
%% method: 无
%% score: 4
%% degree: 650
%% duplicateId: 0
%% body:
一个标有“$220\UV$ $60\UW$”的白炽灯泡，加上的电压 $U$ 由零逐渐增大到 $220\UV$。在此过程中，电压 $U$ 和电流 $I$ 的关系可用图线表示。题中给出的四个图线中，肯定不符合实际的是\xzanswer{ACD}
\fourchoices[answer=ACD, ispicture=true, h=2.8cm]
{\includesvg{figs/17a}}
{\includesvg{figs/17b}}
{\includesvg{figs/17c}}
{\includesvg{figs/17d}}

%% answer: ACD
%% memo:
\memoanswer{
由于金属的电阻率随温度的升高而增大，所以电阻值随电压的增大而增大，即 $U-I$ 曲线的各点的割线的斜率增大，故 B 正确，本题要求选“肯定不符合实际的是”，所以选 ACD。

这里强调说明的是：由电阻的定义式 $R=\frac{U}{I}$ 知，在 $U-I$ 图象中，图象上各点到原点的连线（割线）的斜率等于该点的电阻值，而不是过该点的切线的斜率等于该点的电阻值。如图所示，A 图：电阻一直不变；B 图：$R_{3}>R_{2}>R_{1}$，所以电阻一直增大；C 图：$R_{3}<R_{2}<R_{1}$，所以电阻一直减小；D 图：$R_{2}>R_{1}$，$R_{3}<R_{2}$，所以电阻先增大后减小。

\onepicture[w=8cm]{figs/17_an.png}

而由 $v=\frac{\Delta x}{\Delta t}$ 知，某点的速度 $v$ 是 $x-t$ 图象中的该点的切线的斜率，二者容易混淆。
}

\item
%% number: 18
%% paperName: 1993年全国高考第 18 题 4 分
%% typeId: 2
%% type: 多选题
%% chapter: 运动和力
%% point: 动量守恒定律
%% method: 无
%% score: 4
%% degree: 650
%% duplicateId: 0
%% body:
在质量为 $M$ 的小车中挂有一单摆，摆球的质量为 $m_{0}$。小车（和单摆）以恒定的速度 $V$ 沿光滑水平地面运动，与位于正对面的质量为 $m$ 的静止木块发生碰撞，碰撞的时间极短。在此碰撞过程中，下列哪个或哪些说法是可能发生的？\xzanswer{BC}
\fourchoices[answer=BC]
{小车、木块、摆球的速度都发生变化，分别变为 $v_{1}$、$v_{2}$、$v_{3}$，满足（$M+m_{0}$）$V=Mv_{1}+mv_{2}+m_{0}v_{3}$}
{摆球的速度不变，小车和木块的速度变为 $v_{1}$ 和 $v_{2}$，满足 $MV=Mv_{1}+mv_{2}$}
{摆球的速度不变，小车和木块的速度都变为 $v$，满足 $MV=$（$M+m$）$v$}
{小车和摆球的速度都变为 $v_{1}$，木块的速度变为 $v_{2}$，满足（$M+m_{0}$）$V=$（$M+m_{0}$）$v_{1}+mv_{2}$}

%% answer: BC
%% memo:
\memoanswer{
两车碰撞时，由于摆球不受冲力（时间极短），所以速度不变，故 AD 错误；两车之间碰撞，B 为一般结果，C 为完全非弹性碰撞的结果，都符合动量守恒，故 BC 正确。
}

\item
%% number: 19
%% paperName: 1993年全国高考第 19 题 4 分
%% typeId: 2
%% type: 多选题
%% chapter: 电场
%% point: 带电粒子在交变电场中的运动
%% method: 无
%% score: 4
%% degree: 450
%% duplicateId: 0
%% body:
图中 A、B 是一对中间开有小孔的平行金属板，两小孔的连线与金属板面相垂直，两极板的距离为 $l$。两极板间加上低频交流电压，A 板电势为零，B 板电势 $u=U_{0}\cos\omega t$。现有一电子在 $t=0$ 时穿过 A 板上的小孔射入电场。设初速度和重力的影响均可忽略不计。则电子在两极板间可能\xzanswer{AC}
\onepicture[w=3cm]{figs/19.png}
\fourchoices[answer=AC]
{以 AB 间的某一点为平衡位置来回振动}
{时而向 B 板运动，时而向 A 板运动，但最后穿出 B 板}
{一直向 B 板运动，最后穿出 B 板，如果 $\omega$ 小于某个值 $\omega_{0}$，$l$ 小于某个值 $l_{0}$}
{一直向 B 板运动，最后穿出 B 板，而不论 $\omega$、$l$ 为任何值}

%% answer: AC
%% memo:
\memoanswer{
B 板电势（即 B、A 板间电压）的波形图如图所示：

\onepicture[w=6cm]{figs/19_an.png}

在 $0\sim1$ 时间内，B 板为正向电压，电子从 A 板小孔加速向 B 板运动；$1\sim2$ 时间内，减速向 B 板运动，至 2 时刻速度为 0；$2\sim3$ 时间内，反向加速运动；$3\sim4$ 时间内，反向减速运动，至此为一周期，故电子以 2 时刻所到达的位置为平衡位置来回振动，故 A 正确、D 错误；设 $t=2$ 时刻所到达的位置距 A 板为 $l_{0}$，若 $l<l_{0}$，则电子穿出 B 板，故 C 亦有可能，故 B 错误、C 正确。
}

\item
%% number: 20
%% paperName: 1993年全国高考第 20 题 3 分
%% typeId: 3
%% type: 填空题
%% chapter: 原子和原子核
%% point: 玻尔模型
%% method: 无
%% score: 3
%% degree: 850
%% duplicateId: 0
%% body:
用电磁波照射某原子，使它从能量为 $E_{1}$ 的基态跃迁到能量为 $E_{2}$ 的激发态，该电磁波的频率等于 \tkanswer[2.5cm]{$\frac{E_{2}-E_{1}}{h}$}。

%% answer: $\frac{E_{2}-E_{1}}{h}$
\jdanswer{
$\frac{E_{2}-E_{1}}{h}$
}

%% memo:
\memoanswer{
由玻尔理论得 $E_{2}-E_{1}=h\nu$，解得 $\nu=\frac{E_{2}-E_{1}}{h}$。
}

\item
%% number: 21
%% paperName: 1993年全国高考第 21 题 3 分
%% typeId: 3
%% type: 填空题
%% chapter: 力物体的平衡
%% point: 三力平衡
%% method: 无
%% score: 3
%% degree: 650
%% duplicateId: 0
%% body:
两根长度相等的轻绳，下端悬挂一质量为 $m$ 的物体，上端分别固定在水平天花板上的 M、N 点，M、N 两点间的距离为 $s$，如图所示。已知两绳所能经受的最大拉力均为 $T$，则每根绳的长度不得短于 \tkanswer[3cm]{$\frac{Ts}{\sqrt{4T^{2}-m^{2}g^{2}}}$}。
\onepicture[align=right, w=5cm]{\includesvg{figs/21}}

%% answer: $\frac{Ts}{\sqrt{4T^{2}-m^{2}g^{2}}}$
\jdanswer{
$\frac{Ts}{\sqrt{4T^{2}-m^{2}g^{2}}}$
}

%% memo:
\memoanswer{
根据力的合成和分解的平行四边形法则，受力分析，作图如下。

\onepicture[w=4.5cm]{figs/21_an.png}

有 $\cos\alpha=\frac{mg}{2T}$，$\sin\alpha=\frac{s}{2L}$，$L$ 为绳长，即 OM，

联立解得 $L=\frac{Ts}{\sqrt{4T^{2}-m^{2}g^{2}}}$。
}

\item
%% number: 22
%% paperName: 1993年全国高考第 22 题 3 分
%% typeId: 3
%% type: 填空题
%% chapter: 直线运动
%% point: 游标卡尺和螺旋测微仪
%% method: 无
%% score: 3
%% degree: 850
%% duplicateId: 0
%% body:
有一游标卡尺，主尺的最小分度是 $1\Umm$，游标上有 20 个小的等分刻度。用它测量一工件的长度，如图所示，图示的读数是 \tkanswer[2cm]{104.05}$\Umm$。
\onepicture[align=right, w=7cm]{figs/22.png}

%% answer: 104.05 毫米
\jdanswer{
104.05
}

%% memo:
\memoanswer{
主尺读 $104\Umm$，游标尺读数为 $1\times0.05\Umm=0.05\Umm$，故读数为 $104\Umm+1\times0.05\Umm=104.05\Umm$。
}

\item
%% number: 23
%% paperName: 1993年全国高考第 23 题 3 分
%% typeId: 4
%% type: 实验题
%% chapter: 机械振动
%% point: 单摆测重力加速度
%% method: 无
%% score: 3
%% degree: 650
%% duplicateId: 0
%% body:
一位同学用单摆做测重力加速度的实验。他将摆挂起后，进行了如下操作：

（A）测摆长 $L$：用米尺量出摆线的长度。

（B）测周期 $T$：将摆球拉起，然后放开。在摆球某次通过最低点时，按下秒表开始计时，同时将此次通过最低点作为第一次，接着一直数到摆球第 60 次通过最低点时，按下秒表停止计时。读出这段时间 $t$，算出单摆的周期 $T=\frac{t}{60}$。

（C）将测得的 $L$ 和 $T$ 代入单摆的周期公式 $T=2\pi\sqrt{\frac{l}{g}}$，算出 $g$，将它作为实验的最后结果写入实验报告中。

指出上面步骤中遗漏或错误的地方，写出该步骤的字母，并加以改正（不要求进行误差分析）。

%% answer: （A）应量出悬点到摆球球心之间的长度为摆长；（B）$T=\frac{t}{29.5}$；（C）$g$ 应测量多次，然后取平均值才能作为实验的最后结果。
\jdanswer{
（A）错，应量出悬点到摆球球心之间的长度为摆长；

（B）错，$T=\frac{t}{29.5}$；

（C）错，$g$ 应测量多次，然后取平均值才能作为实验的最后结果。
}

%% memo:
\memoanswer{
（A）错，应量出悬点到摆球球心之间的长度为摆长；

（B）错，$T=\frac{t}{29.5}$；

（C）错，$g$ 应测量多次，然后取平均值才能作为实验的最后结果。
}

\item
%% number: 24
%% paperName: 1993年全国高考第 24 题 3 分
%% typeId: 3
%% type: 填空题
%% chapter: 运动和力
%% point: 动量和能量
%% method: 无
%% score: 3
%% degree: 450
%% duplicateId: 0
%% body:
如图所示，A、B 是位于水平桌面上的两个质量相等的小木块，离墙壁的距离分别为 $L$ 和 $l$，与桌面之间的滑动摩擦系数分别为 $\mu_{A}$ 和 $\mu_{B}$。今给 A 以某一初速度，使之从桌面的右端向左运动。假定 A、B 之间，B 与墙之间的碰撞时间都很短，且碰撞中总动能无损失。若要使木块 A 最后不从桌面上掉下来，则 A 的初速度最大不能超过 \tkanswer[3.5cm]{$\sqrt{4g\left[\mu_{A}(L-l)+\mu_{B}l\right]}$}。
\onepicture[align=right, w=5.5cm]{figs/24.png}

%% answer: $\sqrt{4g\left[\mu_{A}(L-l)+\mu_{B}l\right]}$
\jdanswer{
$\sqrt{4g\left[\mu_{A}(L-l)+\mu_{B}l\right]}$
}

%% memo:
\memoanswer{
A 与 B 碰撞及 B 与墙碰撞都没有机械能损失，且 A 与 B 质量相等，故发生碰撞时，速度发生交换；A 与 B 在运动过程中，只有摩擦力做功损失动能，由动能定理得

$2\mu_{A}mg(L-l)+2\mu_{B}mgl\geq\frac{1}{2}mv_{0}^{2}$，

解得 $v_{0}\leq\sqrt{4g[\mu_{A}(L-l)+\mu_{B}l]}$。
}

\item
%% number: 25
%% paperName: 1993年全国高考第 25 题 3 分
%% typeId: 7
%% type: 作图题
%% chapter: 电路
%% point: 电路中的图象
%% method: 无
%% score: 3
%% degree: 650
%% duplicateId: 0
%% body:
如图 A 所示的电路中，两二极管均可视为理想二极管，$R_{1}=R_{2}$。a 端对 b 端的电压与时间的关系如图 B 的上图所示。请在图 B 的下图中作出 a 端对 c 点的电压与时间的关系图线（最少画一个周期，可用铅笔作图）。
\onepicture[w=9cm]{figs/25.png}

%% answer: 如图
\jdanswer{
如图
}

%% memo:
\memoanswer{
作图时应使图线的上半部和下半部看起来基本像是峰值不同的正弦曲线的一部分，图线的最高点、最低点及与横轴的交点位置必须正确。

当 a 正 b 负时，电流经 $R_{1}$、二极管（上）、$R_{2}$，由于 $R_{1}$、$R_{2}$ 分压，所以 $u_{ac}=\frac{1}{2}u_{ab}$；当 a 负 b 正时，电流经二极管（下）和 $R_{1}$，所以 $u_{ca}=u_{ba}$，所以图象如图所示。

\onepicture[w=8cm]{figs/25_an.png}
}

\item
%% number: 26
%% paperName: 1993年全国高考第 26 题 4 分
%% typeId: 4
%% type: 实验题
%% chapter: 电路
%% point: 电表改装
%% method: 无
%% score: 4
%% degree: 450
%% duplicateId: 0
%% body:
将量程为 $100\UuA$ 的电流表改装成量程为 $1\UmA$ 的电流表，并用一标准电流表与改装后的电流表串联，对它进行校准（核对）。改装及校准所用器材的实物图如下（其中标准电流表事先已与一固定电阻串联，以防烧表）。校准时要求通过电流表的电流能从 $0$ 连续调到 $1\UmA$。试按实验要求在所给的实物图上连线。
\onepicture[w=9cm]{figs/26.png}

%% answer: 如图
\jdanswer{
如图
}

%% memo:
\memoanswer{
电阻箱电阻必须与微安表并联，滑动变阻器要接成分压电路才能满足要求：电流从 $0$ 到 $1\UmA$，若接成限流电路，电流不能到达 $0$ 及附近。

\onepicture[w=9cm]{figs/26_an.png}
}

\item
%% number: 27
%% paperName: 1993年全国高考第 27 题 4 分
%% typeId: 3
%% type: 填空题
%% chapter: 光学
%% point: 透镜
%% method: 无
%% score: 4
%% degree: 450
%% duplicateId: 0
%% body:
某人透过焦距为 $10\Ucm$，直径为 $4.0\Ucm$ 的薄凸透镜观看方格纸，每个方格的边长均为 $0.30\Ucm$。他使透镜的主轴与方格纸垂直，透镜与纸面相距 $10\Ucm$，眼睛位于透镜主轴上离透镜 $5.0\Ucm$ 处。问他至多能看到同一行上几个完整的方格？答：\tkanswer[1.5cm]{26}。

%% answer: 26
\jdanswer{
26
}

%% memo:
\memoanswer{
如图所示，画出经凸透镜折射进入眼睛 C 的两条边界光线 $A'C$ 和 $B'C$。作平行于 $A'C$ 的过光心 O 的光线交焦平面于 A，则 A 为进入人眼光线的一个边界点，作 $AF=FB$，连 $BB'$、$BC$，则 AB 为人眼能透过透镜看到的最大范围，由图可知 $\triangle AFO\sim\triangle A'OC$，则有

\onepicture[w=6cm]{figs/27_an.png}

$\frac{AF}{A'O}=\frac{FO}{OC}$

解得 $AF=\frac{FO}{OC}A'O=4\Ucm$

故 $AB=2AF=8\Ucm$

观察到的格数为 $n=\frac{AB}{0.3}=26.7$，因此观察到完整的方格数 26 个。
}

\item
%% number: 28
%% paperName: 1993年全国高考第 28 题 5 分
%% typeId: 6
%% type: 计算题
%% chapter: 气体性质
%% point: 气体连接体
%% method: 无
%% score: 5
%% degree: 550
%% duplicateId: 0
%% body:
一个密闭的气缸，被活塞分成体积相等的左右两室，气缸壁与活塞是不导热的，它们之间没有摩擦。两室中气体的温度相等，如图所示。现利用右室中的电热丝对右室中的气体加热一段时间。达到平衡后，左室的体积变为原来体积的 $\frac{3}{4}$，气体温度 $T_{1}=300\UK$，求右室气体的温度。
\onepicture[align=right, w=6cm]{figs/28.png}

%% answer: $500\UK$
\jdanswer{
$500\UK$
}

%% memo:
\memoanswer{
设加热前，左室气体的体积为 $V_{0}$，温度为 $T_{0}$，压强为 $p_{0}$。加热后，气体的体积为 $V_{1}$，温度为 $T_{1}$，压强为 $p_{1}$，则有：

$\frac{p_{0}V_{0}}{T_{0}}=\frac{p_{1}V_{1}}{T_{1}}$ ①

由题意知，加热前右室气体的体积、压强和温度也分别为 $V_{0}$、$p_{0}$ 和 $T_{0}$，若加热后变为 $V_{2}$、$p_{2}$ 和 $T_{2}$，则有

$\frac{p_{0}V_{0}}{T_{0}}=\frac{p_{2}V_{2}}{T_{2}}$ ②

由题意知

$p_{1}=p_{2}$，$V_{1}=\frac{3}{4}V_{0}$，$V_{2}=V_{0}+\frac{1}{4}V_{0}=\frac{5}{4}V_{0}$。

代入①②式解得

$T_{2}=\frac{5}{3}T_{1}=\frac{5}{3}\times300\UK=500\UK$。

评分标准：全题 5 分。列出左、右两室气体的气态方程占 1 分；加热前和加热后，两室中气体的压强都相等各占 1 分；求得加热后右室气体的体积占 1 分；求得最后结果占 1 分。
}

\item
%% number: 29
%% paperName: 1993年全国高考第 29 题 5 分
%% typeId: 6
%% type: 计算题
%% chapter: 电磁感应
%% point: 电磁感应受力分析
%% method: 无
%% score: 5
%% degree: 550
%% duplicateId: 0
%% body:
两金属杆 ab 和 cd 长均为 $l$，电阻均为 $R$，质量分别为 $M$ 和 $m$，$M>m$。用两根质量和电阻均可忽略的不可伸长的柔软导线将它们连成闭合回路，并悬挂在水平、光滑、不导电的圆棒两侧。两金属杆都处在水平位置，如图所示。整个装置处在一与回路平面相垂直的匀强磁场中，磁感应强度为 $B$。若金属杆 ab 正好匀速向下运动，求运动的速度。
\onepicture[align=right, w=5cm]{\includesvg{figs/29}}

%% answer: $v=\frac{(M-m)gR}{2B^{2}l^{2}}$
\jdanswer{
$v=\frac{(M-m)gR}{2B^{2}l^{2}}$
}

%% memo:
\memoanswer{
设磁场方向垂直纸面向里，由法拉第电磁感应定律得，

ab 中的感应电动势为 $E_{1}=Blv$，方向由 $a\rightarrow b$。

cd 中的感应电动势 $E_{2}=Blv$，方向由 $d\rightarrow c$。

回路中电流方向由 $a\rightarrow b\rightarrow d\rightarrow c$，大小为

$I=\frac{E_{1}+E_{2}}{2R}=\frac{2Blv}{2R}=\frac{Blv}{R}$

ab 受到的安培力向上，cd 受到的安培力向下，大小均为

$F_{\text{安}}=BIl=\frac{B^{2}l^{2}v}{R}$

当 ab 匀速下滑时，对 ab 有 $T+F_{\text{安}}=Mg$

对 cd 有 $T=F_{\text{安}}+mg$

式中 $T$ 为杆所受到的导线的拉力

解得 $2F_{\text{安}}=(M-m)g$

$\frac{2B^{2}l^{2}v}{R}=(M-m)g$

解得 $v=\frac{(M-m)gR}{2B^{2}l^{2}}$

评分标准：全题 5 分。正确求得电流 $I$ 值，得 2 分；求得作用于两杆的安培力得 1 分；求得两杆做匀速运动时力的平衡式得 1 分；求得速度再得 1 分。

若设磁场方向垂直纸面向外，正确的，同样给分。
}

\item
%% number: 30
%% paperName: 1993年全国高考第 30 题 6 分
%% typeId: 6
%% type: 计算题
%% chapter: 力物体的平衡
%% point: 力矩平衡
%% method: 无
%% score: 6
%% degree: 450
%% duplicateId: 0
%% body:
有一准确的杆秤。今只给你一把有刻度的直尺，要求用它测出这杆秤的秤砣的质量。试导出表示秤砣质量的公式，并说明所需测量的量。

%% answer: $m=\frac{Md}{l_{0}+l}$；从秤杆上读出最大秤量 $M$，用直尺测出 $d$ 和从 A 点到 D 点的距离（$l_{0}+l$），代入上式即可求得 $m$。
\jdanswer{
$m=\frac{Md}{l_{0}+l}$

从秤杆上读出最大秤量 $M$，用直尺测出 $d$ 和从 A 点到 D 点的距离（$l_{0}+l$），代入上式即可求得 $m$。
}

%% memo:
\memoanswer{
解法一：

秤的结构如图所示，秤钩 B 到提钮的距离为 $d$，零刻度（即定盘星）A 到提钮的距离为 $l_{0}$，满刻度 D 到提钮的距离为 $l$，秤杆和秤钩所受的重力为 $P$，秤水平时，$P$ 对提钮的力臂为 $d_{0}$。设秤砣的质量为 $m$，杆秤的最大秤量为 $M$。

\onepicture[w=7cm]{figs/30_an1.jpg}

当空秤平衡时，有

$mgl_{0}=Pd_{0}$ ①

当满秤量平衡时，有

$Mgd=Pd_{0}+mgl$ ②

解①、②两式得

$m=\frac{Md}{l_{0}+l}$ ③

从秤杆上读出最大秤量 $M$，用直尺测出 $d$ 和从 A 点到 D 点的距离（$l_{0}+l$），代入上式即可求得 $m$。

评分标准：全题 6 分。①、②两式都正确给 3 分，只有一式正确给 1 分；求得③式再给 1 分；说出用直尺测量 $d$、$(l_{0}+l)$ 两个量给 2 分，缺少其中任何一个量都不给这 2 分；说分别测量 $d$、$l_{0}$、$l$ 的也给这 2 分，但缺少其中任何一个量都不给这 2 分。把定盘星放在提钮的另一侧，正确的，同样给分。

解法二：

秤的结构如图所示。设秤钩 B 到提钮的距离为 $d$，秤杆和秤钩所受的重力为 $P$。秤水平时，$P$ 对提钮的力臂为 $d_{0}$，秤砣的质量为 $m$。

\onepicture[w=7cm]{figs/30_an2.jpg}

设想先把秤砣挂在秤杆读数为 $M_{1}$ 处，该处到提钮的距离为 $l_{1}$，平衡时有：

$M_{1}gd=Pd_{0}+mgl_{1}$ ①

再把秤砣挂在秤杆的读数为 $M_{2}$ 处，该处到提钮的距离为 $l_{2}$，平衡时有：

$M_{2}gd=Pd_{0}+mgl_{2}$ ②

解①、②两式得

$m=\frac{(M_{2}-M_{1})d}{l_{2}-l_{1}}$ ③

从秤杆上读出 $M_{1}$、$M_{2}$，用直尺测得 $d$ 和从 $M_{1}$ 处到 $M_{2}$ 处的距离 $l_{2}-l_{1}$，代入上式即得 $m$。评分标准：与解法一相同。

解法三：

秤的结构如图所示，秤钩 B 到提钮的距离为 $d$，A 是零刻度（即定盘星），D 是满刻度。设秤砣的质量为 $m$。当把秤砣挂放在零刻度上，秤平衡时秤钩是空的。若把秤砣从 A 点移到 D 点，对提钮增加的力矩为 $mgl$，$l$ 为 AD 间的距离，则在秤钩上挂一质量为 $M$ 的物体后，秤又平衡。这表示重物对提钮增加的力矩 $Mgd$ 与 $mgl$ 大小相等

\onepicture[w=7cm]{figs/30_an3.jpg}

即

$Mgd=mgl$ ①

解得 $m=\frac{Md}{l}$

从秤上读出最大秤量 $M$，用直尺量出 $d$ 和 $l$，代入上式即求出 $m$。

评分标准：全题 6 分。在分析正确，说理清楚的前提下，直接得到①式给 4 分；说出用直尺测量 $l$、$d$ 两个量给 2 分，缺少其中任何一个量，不给这 2 分。
}

\item
%% number: 31
%% paperName: 1993年全国高考第 31 题 8 分
%% typeId: 6
%% type: 计算题
%% chapter: 运动和力
%% point: 牛顿定律的应用
%% method: 无
%% score: 8
%% degree: 250
%% duplicateId: 0
%% body:
一平板车，质量 $M=100\Ukg$，停在水平路面上，车身的平板离地面的高度 $h=1.25\Um$，一质量 $m=50\Ukg$ 的小物块置于车的平板上，它到车尾端的距离 $b=1.00\Um$，与车板间的滑动摩擦系数 $\mu=0.20$，如图所示。今对平板车施一水平方向的恒力，使车向前行驶，结果物块从车板上滑落。物块刚离开车板的时刻，车向前行驶的距离 $s_{0}=2.0\Um$。求物块落地时，落地点到车尾的水平距离 $s$。不计路面与平板车间以及轮轴之间的摩擦。
\onepicture[align=right, w=6cm]{figs/31.png}

%% answer: $s=1.6\Um$
\jdanswer{
$s=1.6\Um$
}

%% memo:
\memoanswer{
解法一：

设作用于平板车的水平恒力为 $F$，物块与车板间的摩擦力为 $f=\mu mg$，自车启动至物块开始离开车板经历的时间为 $t$，物块开始离开车板时的速度为 $v$，车的速度为 $V$。系统的运动过程如图所示。

\onepicture[w=8cm]{figs/31_an.jpg}

由动能定理得

$(F-f)s_{0}=\frac{1}{2}MV^{2}$ ①，

$f(s_{0}-b)=\frac{1}{2}mv^{2}$ ②，

由动量定理得

$(F-f)t=MV$ ③

$ft=mv$ ④

$f=\mu mg$ ⑤

由①②得 $\frac{F-f}{f}\cdot\frac{s_{0}}{s_{0}-b}=\frac{MV^{2}}{mv^{2}}$ ⑥，

由③④式得 $\frac{F-f}{f}=\frac{MV}{mv}$ ⑦，

由②⑤式得 $v=\sqrt{2\mu g(s_{0}-b)}=2\Ums$，

由⑥⑦式得 $V=\frac{s_{0}}{s_{0}-b}v=\frac{2}{2-1}\times2=4\Ums$，

由①式得

$F=f+\frac{1}{2}\frac{MV^{2}}{s_{0}}=\mu mg+\frac{1}{2}\frac{MV^{2}}{s_{0}}=500\UN$，

物块离开车板后作平抛运动，其水平速度 $v$，设经历的时间为 $t_{1}$，所经过的水平距离为 $s_{1}$，则有

$s_{1}=vt_{1}$ ⑧

$h=\frac{1}{2}gt_{1}^{2}$ ⑨，

由⑨式得 $t_{1}=\sqrt{\frac{2h}{g}}=0.5\Us$，$s_{1}=2\times0.5=1\Um$，

物块离开平板车后，若车的加速度为 $a$，由牛顿第二定律得

$a=\frac{F}{M}=\frac{500}{100}=5\Umsq$，

车运动的距离为

$s_{2}=Vt_{1}+\frac{1}{2}at_{1}^{2}=2.652\Um$，取两位，$s_{2}=2.6\Um$。

于是

$s=s_{2}-s_{1}=2.6-1=1.6\Um$

落地点到车尾的水平距离为 $1.6\Um$。

评分标准：全题 8 分。正确求得物块开始离开车板时刻的物块速度 $v$ 给 1 分，车的速度 $V$ 给 2 分；求得作用于车的恒力 $F$ 再给 1 分。正确求得物块离开车板后平板车的加速度给 1 分。正确分析物块离开车板后的运动，并求得有关结果，正确求出物块下落过程中车的运动距离 $s_{2}$ 并由此求得 $s$ 的正确数值，共给 3 分。最后结果有错，不给这 3 分。

解法二：

设作用于平板车的水平恒力为 $F$，物块与车板间的摩擦力为 $f$，自车启动至物块离开车板经历的时间为 $t$，在这过程中，车的加速度为 $a_{1}$，物块的加速度为 $a_{2}$。则有

$F-f=Ma_{1}$ ①

$f=ma_{2}$ ②

$f=\mu mg$ ③

以及

$s_{0}=\frac{1}{2}a_{1}t^{2}$ ④，$s_{0}-b=\frac{1}{2}a_{2}t^{2}$ ⑤，

由②、③两式得

$a_{2}=\mu g=0.2\times10=2\Umsq$

由④、⑤两式得

$a_{1}=\frac{s_{0}}{s_{0}-b}a_{2}=4\Umsq$，

由①、③两式得

$F=\mu mg+Ma_{1}=0.2\times50\times10+100\times4=500\UN$

物块开始离开车板时刻，物块和车的速度分别为 $v$ 和 $V$，则

$V=\sqrt{2a_{1}s_{0}}=4\Ums$，

$v=\sqrt{2a_{2}(s_{0}-b)}=2\Ums$，

物块离车板后作平抛运动，其水平速度为 $v$，所经历的时间为 $t_{1}$，走过的水平距离为 $s_{1}$，则有 $s_{1}=vt_{1}$ ⑥

$h=\frac{1}{2}gt_{1}^{2}$ ⑦

解得 $t_{1}=\sqrt{\frac{2h}{g}}=0.5\Us$，$s_{1}=1\Um$，

在这段时间内车的加速度

$a=\frac{F}{M}=\frac{500}{100}=5\Umsq$，

车运动的距离为

$s_{2}=Vt_{1}+\frac{1}{2}at_{1}^{2}=2.652\Um$，取两位，$s_{2}=2.6\Um$，

故 $s=s_{2}-s_{1}=2.6-1=1.6\Um$

落地点到车尾的水平距离为 $1.6\Um$。

评分标准：全题 8 分。正确求得物块离开车板前，物块和车的加速度 $a_{1}$、$a_{2}$，占 2 分，求得物块开始离开车板时刻的速度 $v$ 和此时车的速度 $V$ 占 1 分，求得作用于车的恒力 $F$ 占 1 分。正确求得物块离开车板后，车的加速度 $a$ 占 1 分。正确分析物块离开车板后物块的运动并求得有关结果，正确求得物块下落过程中车的运动距离，并由此求得 $s$ 的正确结果，共占 3 分。最后结果错误，不给这 3 分。
}

\end{enumerate}

\end{document}
